\chapter{Internship outcomes and skills}
\section{Learning the molecular context}
The internship required more than applying a familiar statistical method to a new table. The assigned programme began with the concepts needed to interpret the data: amino acids, peptide bonds, protein structure and intrinsic disorder. Force fields, molecular dynamics and dihedral angles connected those concepts to simulation. The final task brought them together through examination of a clustering analysis and an independent investigation.

I used GROMACS, VMD and Python during the internship. Learning to use GROMACS was one of the more challenging parts of the internship for me. I found the process demanding at first, but I also enjoyed becoming more familiar with the software.

Saved configurations need not be independent observations, and their frequency in a biased simulation differs from their equilibrium probability. These points connect the molecular setting to statistics while showing why the origin of the data matters.

\section{Turning an observation into an analysis}
The Python project provides a concrete example. The descriptor analysis found a strong relationship between radius of gyration and SASA. Rather than treating this as a reason to remove SASA immediately, the independent experiment compared fitted assignments with and without it. The same fitting rows kept the comparison focused on the feature choice. Repeated initializations exposed variation that a single fit would have hidden.

Small programming decisions also became explicit parts of the analysis. Feature scaling is fitted on a specified sample and reused for later observations. A test checks that a large evaluation value cannot alter that transformation. Original frame identities are retained so that label arrays refer to the same observations. These details determine whether a difference comes from the method under investigation or an unintended change in preprocessing.

\begin{codepanel}{Fit once, reuse the same scale}{RobustCluster-Aβ / Python}
\lstinputlisting[style=editorialcode]{generated/fit_scaler.py}
\end{codepanel}


The modules, tests, configurations, source tables and plots document the work produced. Successful execution remains distinct from establishing the physical meaning of a cluster.

\section{My molecular-software interface projects}
I also created two personal software projects, \href{https://github.com/Alierkn/vmd-mcp}{VMD MCP} and \href{https://github.com/Alierkn/gromacs-mcp}{GROMACS MCP}. They use the Model Context Protocol to expose molecular-software operations through structured inputs and outputs. They are distinct from the Aβ40 descriptor analysis, but provide further examples of connecting scientific programs with Python interfaces.

\par\addvspace{8pt}\noindent
\begin{comparisonpanel}{VMD MCP: frame measurements}
The inspected source constructs Tcl commands for measurement and rendering. Its gyration function passes an atom selection and mass-weighting option to VMD, then returns frame-indexed values.
\end{comparisonpanel}\hfill
\begin{comparisonpanel}{GROMACS MCP: execution interface}
The inspected source includes preparation and analysis helpers. Its background-run interface checks paths and records output in a log.
\end{comparisonpanel}
\par\addvspace{8pt}
Contributor histories remain in the repositories. Appendix G links to the source versions inspected and pairs the code excerpts with genuine screenshots.

These interfaces organize requests; source code alone does not establish that a calculation ran. No reported Aβ40 trajectory is attributed to them. Comparing their outputs with the supplied descriptors would also require compatible selections and units: the inspected VMD gyration function documents angstroms, whereas the table is labelled in nanometres.

\section{What can be carried forward}
The report brings together the assigned learning programme, practical software use and calculations with saved outputs. Distinguishing them gives a clearer account than treating every planned task as completed. GROMACS and VMD use belongs to the practical experience; the descriptor diagnostics and GMM comparison belong to the recorded analysis.

The reusable outcome is a project whose steps can be examined. Each empirical figure has a source table, each calculation has settings, and each unavailable result has a reason. This organization can support further work once the FES units and construction are established. A Leonardo submission template is included for preparation, but no Leonardo job records exist. The reported computations remain local; further simulations or structural analyses will need their own execution records and validation.
